\chapter*{Abstract}
\addcontentsline{toc}{chapter}{Abstract}

Electricity-consumption forecasting is useful only when reported performance reflects genuinely unseen future conditions. This project began with the replication of a published Deep Energy Predictor Model combining a cascaded residual network, a deep neural network, and XGBoost. Although the initial implementation reproduced classification accuracy close to the values announced in the article, a methodological audit showed that the task and validation protocol allowed target information to enter the predictors. The most important mechanisms were direct use of a target-derived variable, the physical relationship between active power, voltage, and current, non-shifted rolling features, preprocessing before data separation, and random splitting of strongly correlated temporal observations. The historical high scores were therefore retained as diagnostic evidence but excluded from forecasting claims.

The problem was reformulated as continuous multi-horizon forecasting with causal windows, chronological train-validation-test separation, train-only fitting of preprocessing operations, inverse transformation, physical-unit metrics, and seasonal baselines. The research then evaluated recurrent, convolutional, patch-based, inverted-Transformer, and global forecasting approaches across IHEPC, Steel Industry, ECL, and Low Carbon London data. Final checkpoint selection is supported by future-time evaluation on households excluded from gradient updates. The 24-hour Global Temporal Fusion Transformer obtains a cold-start macro MAE of 0.1849 kWh per hourly interval, 26.48\% below the seasonal-naive baseline, and outperforms that baseline for 99.0\% of evaluated households. The 168-hour model obtains 0.1973 kWh per hourly interval, a 20.77\% improvement, with 98.40\% of households outperforming the baseline.

These checkpoints were integrated into EnergyAI, a Dockerized Next.js and FastAPI platform with PostgreSQL persistence, authenticated user ownership, CSV ingestion, artifact contracts, forecast history, alerts, actions, analytics export, and explicit readiness checks. Current validation includes a complete 100-test PostgreSQL backend run, successful frontend linting, type checking and production compilation, and a six-project Playwright run with one non-reproduced WebKit reload interruption. A remaining train--serve mismatch is reported explicitly: frozen research metrics use household training-segment scalers, whereas serving uses the latest 336-hour window. The contribution is therefore an evidence-traceable path from flawed evaluation to a controlled operational prototype, not an unsupported claim of deployed accuracy.

\textbf{Keywords:} electricity consumption forecasting; data leakage; time series; Temporal Fusion Transformer; cold-start evaluation; reproducibility; EnergyAI.

\cleardoublepage
\begin{otherlanguage}{french}
\chapter*{Résumé}
\addcontentsline{toc}{chapter}{Résumé}
\enlargethispage{2\baselineskip}

La prévision de la consommation électrique n'est utile que lorsque les performances annoncées reflètent réellement des conditions futures non observées. Ce projet a débuté par la réplication d'un Deep Energy Predictor Model combinant un réseau résiduel en cascade, un réseau neuronal profond et XGBoost. Bien que la première implémentation ait reproduit une exactitude de classification proche des valeurs annoncées dans l'article, un audit méthodologique a montré que la tâche et le protocole de validation permettaient à des informations dérivées de la cible d'entrer dans les variables explicatives. Les principaux mécanismes identifiés sont l'utilisation directe d'une variable liée à la cible, la relation physique entre puissance active, tension et intensité, des fenêtres glissantes non décalées, le prétraitement avant séparation et une division aléatoire d'observations temporelles fortement corrélées. Les scores historiques élevés sont donc conservés comme éléments de diagnostic, mais exclus des contributions de prévision.

Le problème a été reformulé en prévision continue multi-horizon avec fenêtres causales, séparation chronologique entraînement-validation-test, ajustement des prétraitements sur l'entraînement uniquement, transformation inverse, métriques en unités physiques et baselines saisonnières. L'étude a ensuite comparé des approches récurrentes, convolutionnelles, fondées sur des patches, des Transformers inversés et des modèles globaux sur les jeux IHEPC, Steel Industry, ECL et Low Carbon London. La sélection finale des checkpoints repose sur une évaluation temporelle de ménages exclus des mises à jour de gradient. Le Global Temporal Fusion Transformer à 24 heures atteint une MAE macro cold-start de 0,1849 kWh par intervalle horaire, soit 26,48\% de mieux que la baseline saisonnière, et la surpasse pour 99,0\% des ménages évalués. À 168 heures, le modèle atteint 0,1973 kWh par intervalle horaire, soit une amélioration de 20,77\%, et surpasse la baseline pour 98,40\% des ménages.

Ces checkpoints sont intégrés dans EnergyAI, une plateforme Dockerisée fondée sur Next.js, FastAPI et PostgreSQL, avec authentification, isolation des données, import CSV, contrats d'artefacts, historique des prévisions, alertes, actions, export analytique et contrôles explicites de disponibilité. Une divergence entraînement--service demeure explicitement signalée : les métriques de recherche utilisent les statistiques du segment d'entraînement de chaque ménage, tandis que le service utilise la fenêtre récente de 336 heures. La contribution essentielle est donc une chaîne de preuves allant d'une évaluation invalide à un prototype opérationnel contrôlé, et non une affirmation non vérifiée de précision en déploiement.

\textbf{Mots-clés :} prévision de consommation électrique ; fuite de données ; séries temporelles ; Temporal Fusion Transformer ; évaluation cold-start ; reproductibilité ; EnergyAI.
\end{otherlanguage}
